% ৪.১৩ অনুশীলনী
\section{অনুশীলনী}

\begin{learningbox}
এই অনুশীলন শেষে শিক্ষার্থী পারবে:
\begin{itemize}
\item HTML tag, attribute, table, form ও hyperlink নিয়ে board-style প্রশ্ন সমাধান করতে।
\item একটি simple webpage-এর structure লিখতে।
\item website publishing ও design workflow ব্যাখ্যা করতে।
\end{itemize}
\end{learningbox}

\begin{practicebox}
\textbf{MCQ}
\begin{enumerate}
\item HTML-এর পূর্ণরূপ কোনটি?\\
ক. Hyper Transfer Markup Language \quad খ. HyperText Markup Language \quad গ. HighText Machine Language \quad ঘ. Hyperlink Text Mode Language
\item Image দেখানোর tag কোনটি?\\
ক. \texttt{<a>} \quad খ. \texttt{<img>} \quad গ. \texttt{<table>} \quad ঘ. \texttt{<form>}
\item Hyperlink তৈরিতে কোন attribute সবচেয়ে গুরুত্বপূর্ণ?\\
ক. src \quad খ. href \quad গ. alt \quad ঘ. method
\item Table row তৈরির tag কোনটি?\\
ক. \texttt{<tr>} \quad খ. \texttt{<td>} \quad গ. \texttt{<th>} \quad ঘ. \texttt{<br>}
\item Form-এর sensitive data পাঠাতে কোন method বেশি উপযোগী?\\
ক. GET \quad খ. POST \quad গ. LINK \quad ঘ. IMG
\end{enumerate}
\end{practicebox}

\begin{practicebox}
\textbf{সংক্ষিপ্ত প্রশ্ন}
\begin{enumerate}
\item ওয়েব পেজ কী?
\item Static ও dynamic website-এর পার্থক্য লেখ।
\item Tag, element ও attribute উদাহরণসহ ব্যাখ্যা কর।
\item \texttt{alt} attribute কেন দরকার?
\item \texttt{rowspan} ও \texttt{colspan}-এর কাজ লেখ।
\item Domain name ও web hosting-এর পার্থক্য লেখ।
\end{enumerate}
\end{practicebox}

\begin{practicebox}
\textbf{সৃজনশীল অনুশীলন ১}\\
একটি কলেজ website-এ Home, Notice, Routine, Gallery ও Contact page আছে। Notice page-এ PDF link, Routine page-এ table এবং Contact page-এ form রাখা হয়েছে।

\begin{enumerate}
\item ক. Home page কী?
\item খ. Hyperlink কেন website navigation-এর জন্য গুরুত্বপূর্ণ?
\item গ. Routine page-এর জন্য table tag ব্যবহার করে structure ব্যাখ্যা কর।
\item ঘ. Contact page-এর form design করতে কোন কোন tag ও validation দরকার হবে বিশ্লেষণ কর।
\end{enumerate}
\end{practicebox}

\begin{practicebox}
\textbf{সৃজনশীল অনুশীলন ২}\\
একজন শিক্ষার্থী একটি personal portfolio তৈরি করছে। সে একটি profile image, education table, project link এবং feedback form রাখতে চায়। websiteটি mobile-এও সুন্দর দেখাতে হবে।

\begin{enumerate}
\item ক. Responsive design কী?
\item খ. Image-এর \texttt{alt} attribute কেন ব্যবহার করা হয়?
\item গ. portfolio page-এর একটি layout plan তৈরি কর।
\item ঘ. website publish করার আগে কী কী test করা উচিত? যুক্তিসহ লেখ।
\end{enumerate}
\end{practicebox}

\begin{recapbox}
\begin{itemize}
\item HTML structure দেয়, CSS style দেয়, JavaScript behavior যোগ করে।
\item \texttt{href} link-এর গন্তব্য, \texttt{src} image/file path, \texttt{alt} বিকল্প text।
\item Table: \texttt{table} $\rightarrow$ \texttt{tr} $\rightarrow$ \texttt{th}/\texttt{td}।
\item Form: input গ্রহণ করে; GET search/filter-এ, POST sensitive submission-এ বেশি উপযোগী।
\item Publish করার আগে link, image, form, mobile view ও browser compatibility test করতে হয়।
\end{itemize}
\end{recapbox}

\begin{sourcebox}
Chapter 04 handnote/main sheet, HSC ICT chapter outline এবং standard HTML/web design terminology অনুসরণ করে এই অংশগুলো নতুনভাবে রচিত হয়েছে; কোনো বইয়ের দীর্ঘ অংশ হুবহু ব্যবহার করা হয়নি।
\end{sourcebox}
